\documentclass[12pt]{article}
\usepackage[margin=1in]{geometry}
\usepackage{enumitem}
\usepackage{titlesec}
\usepackage{parskip}
\usepackage{amsmath}
\usepackage{graphicx}
\usepackage{hyperref}
\usepackage{fontenc}

\title{\textbf{Olson 1965 RG Notes}\\
\large Olson, Mancur. \textit{The Logic of Collective Action: Public Goods and the Theory of Groups}.\\
Boston: Harvard University Press, 1965.}
\author{}
\date{}

\begin{document}

\maketitle

\section{Main Idea}



%--------------------------------------------------
\section{General Notes}
%--------------------------------------------------

\begin{itemize}[leftmargin=*, itemsep=4pt]
    \item Collective action problem exists -- rational actors will not support a collective good if they have an incentive to deviate / shirk
    \item Group size as a key variable (the larger the group, the bigger the free-riding problem)
    \item Effective groups often instead provide some degree of private services bundled with common resources, such that there is some distinct individual incentive to contribute. Collective ends should be by-products of operations
    \item Public goods are non-excludable (cannot prohibit access for non-contributors) and non-rival (one person's consumption does not reduce availability)
    \item A group is ``privileged'' if it is rational for at least one member to bear the cost alone (still can introduce a coordination problem though)
    \item Must provide some form of selective incentives where receipt (positive or negative) can be targeted based on contribution
\end{itemize}


\section{Important Specific Factual Claims}

\begin{itemize}[leftmargin=*, itemsep=4pt]
    \item ``Exploitation of the great by the small'' -- those with the largest stakes will usually endure the largest costs (disproportionately) as they are most incentivized to maintain provision
\end{itemize}


\section{Connections}

\begin{itemize}[leftmargin=*, itemsep=4pt]
    \item
\end{itemize}


\section{My Critiques}

\begin{itemize}[leftmargin=*, itemsep=4pt]
    \item
\end{itemize}


\end{document}
